\documentclass[12pt, a4paper]{article}

\usepackage[utf8]{inputenc}

\usepackage[T2A]{fontenc}

\usepackage[left=2cm, right=2cm, top=2cm, bottom=2cm]{geometry}



\begin{document}



\textbf{Вопрос №5.} Что такое множество нулевой меры на прямой, а на плоскости? Привести примеры. Говорят, что число $x$, лежащее в $\mathbb{R}$, является алгебраическим, если оно есть решение полиномиального уравнения с целыми коэффициентами. Показать, что множество алгебраических чисел бесконечно, но имеет нулевую меру.



\vspace{2em}



\textbf{Ответ.} Говорят, что множество $A \subset \mathbb{R}^d$ имеет нулевую меру, если его можно накрыть набором интервалов $\bigcup_{k=1}^\infty I_k$, где длина каждого $I_k$ равна $\frac{\epsilon}{2^k}$ для любого $\epsilon > 0$. Например, на прямой множества нулевой меры: целые числа $\mathbb{Z}$, рациональные числа $\mathbb{Q}$, множество из одной точки.



\vspace{2em}



Для алгебраических чисел рассмотрим полином $a_n x^n + \dots + a_1 x + a_0 = 0$, где $a_i$ — целые коэффициенты. Каждому полиному соответствует конечное число корней, а объединение всех корней всех полиномов счетно. Поскольку счетное объединение счетных множеств остается счетным, множество алгебраических чисел имеет нулевую меру.

\end{document}
